\section*{Problems}

\subsection{Problem Statements}

\subsubsection*{1. Fundamental Level (Simple / Direct)}

\begin{problema}
	\label{prob-ic-1}
	A random sample of $n = 36$ students obtained a sample mean score of $\bar{X} = 78$ points on a mathematics exam. Assume that the population standard deviation is known and exactly equal to $\sigma = 12$ points.
	\begin{enumerate}
		\item Construct a 95\% confidence interval for the population mean $\mu$.
		\item Construct a 99\% confidence interval for the population mean $\mu$.
		\item Which interval has a larger width? Explain the theoretical reason why the width varies as a function of the confidence level.
	\end{enumerate}
\end{problema}

\begin{problema}
	\label{prob-ic-2}
	An LED lightbulb manufacturer wishes to estimate the average lifespan of its new product line. A random sample of $n = 16$ lightbulbs is taken, recording a sample mean of $\bar{X} = 1,200$ hours with a sample standard deviation $s = 80$ hours.
	\begin{enumerate}
		\item Should you use the normal $Z$ distribution or Student's $t$-distribution to compute the confidence interval? Justify your answer.
		\item Construct a 95\% confidence interval for the true average lifespan.
		\item Formally interpret the resulting interval in the context of the problem.
	\end{enumerate}
\end{problema}

\begin{problema}
	\label{prob-ic-fund-new}
	A quality control analyst plans to construct a 98\% parametric confidence interval for the mean tensile strength of a synthetic thread, assuming a known population standard deviation $\sigma$ and a large sample size ($n \ge 30$).
	\begin{enumerate}
		\item Determine the excluded area proportion $\alpha$ and the corresponding area in each extreme tail ($\alpha / 2$).
		\item Calculate the exact critical value $Z_{\alpha/2}$ in the standard normal distribution for this 98\% confidence level.
	\end{enumerate}
\end{problema}

\subsubsection*{2. Operational Level (Intermediate / Developmental)}

\begin{problema}
	\label{prob-ic-3}
	An academic researcher investigates the average weekly time university students dedicate to autonomous study. They select a simple random sample of $n = 25$ students and obtain $\bar{X} = 15.0$ hours along with a sample standard deviation $s = 4.0$ hours.
	\begin{enumerate}
		\item Construct the 90\% confidence interval for the population mean $\mu$.
		\item Construct the 95\% confidence interval for the same mean.
		\item If the researcher wishes to obtain a higher-precision estimate (a narrower confidence interval) while maintaining the 95\% confidence level, what specific methodological modifications should they apply?
	\end{enumerate}
\end{problema}

\begin{problema}
	\label{prob-ic-4}
	A telecommunications firm needs to determine the true average duration of phone calls made by its corporate subscribers. From an audit of $n = 100$ calls, it finds a sample mean $\bar{X} = 4.5$ minutes with $s = 2.0$ minutes.
	\begin{enumerate}
		\item Construct the 95\% confidence interval for the mean duration $\mu$.
		\item Determine the magnitude of the current estimation margin of error ($E$).
		\item If technical management requires reducing the maximum tolerated margin of error to $E = 0.2$ minutes at the same 95\% confidence level, derive and calculate the minimum total sample size $n$ that must be collected.
	\end{enumerate}
\end{problema}

\begin{problema}
	\label{prob-ic-6}
	The engineering department of a battery plant claims in its advertising that the lifespan of its storage cells exceeds 500 recharge cycles. An independent certification laboratory tests $n = 25$ units and obtains $\bar{X} = 490$ cycles with a sample standard deviation $s = 30$ cycles.
	\begin{enumerate}
		\item Construct the 95\% confidence interval for the true average number of cycles $\mu$.
		\item Does the constructed interval contain or exclude the 500-cycle threshold claimed in the advertising?
		\item Explain the technical conclusion regarding the truthfulness of the claim and detail how this interval relates to the outcome of a two-tailed hypothesis test at level $\alpha = 0.05$.
	\end{enumerate}
\end{problema}

\subsubsection*{3. Analytical Level (Advanced / Connection)}

\begin{problema}
	\label{prob-ic-5}
	A comparative clinical nutrition study evaluates the average daily caloric intake between two adult subpopulations (Men and Women). Two independent samples are drawn, yielding the following summary statistics:
	\begin{itemize}
		\item Sample of men: $n_1 = 40$, $\bar{X}_1 = 2,500$ kcal, $s_1 = 300$ kcal.
		\item Sample of women: $n_2 = 45$, $\bar{X}_2 = 2,100$ kcal, $s_2 = 250$ kcal.
	\end{itemize}
	\begin{enumerate}
		\item Calculate the estimated standard error for the difference of independent means $\text{SE}(\bar{X}_1 - \bar{X}_2)$ without assuming that the population variances are necessarily equal.
		\item Construct a 95\% confidence interval for the true difference between the population means $\mu_1 - \mu_2$.
		\item Based on whether the interval includes zero, is there a statistically significant biological or behavioral difference between both genders in their caloric intake?
	\end{enumerate}
\end{problema}

\begin{problema}
	\label{prob-ic-analit-pooled}
	\textbf{Confidence interval for the difference of means with pooled common variance ($S_p^2$):}
	An industrial chemist tests two alternative catalysts ($A$ and $B$) to improve the percentage yield of a thermochemical reaction in small batches. Under the reasonable assumption of normality and equal population variances ($\sigma_1^2 = \sigma_2^2$), the following measurements are obtained:
	\begin{itemize}
		\item Catalyst A: $n_1 = 10$, $\bar{X}_1 = 85.0\%$, $s_1 = 4.0\%$.
		\item Catalyst B: $n_2 = 12$, $\bar{X}_2 = 81.0\%$, $s_2 = 5.0\%$.
	\end{itemize}
	\begin{enumerate}
		\item Derive and calculate the degrees-of-freedom-weighted combined variance (pooled variance) $S_p^2$.
		\item Construct a 95\% confidence interval for the true difference in mean yield $\mu_A - \mu_B$ using the $t$-distribution with $\nu = n_1 + n_2 - 2$ degrees of freedom.
	\end{enumerate}
\end{problema}

\subsubsection*{4. Challenging Level (Experts / Deep Deduction)}

\begin{problema}
	\label{prob-ic-desaf-muestra-bern}
	\textbf{Analytical derivation of the conservative sample size bound for a Bernoulli proportion:}
	When designing an opinion poll or survey to estimate the true population proportion $p \in (0,1)$ with a maximum allowable margin of error $E$ and confidence $(1-\alpha)$, the algebraic expression defining the half-width of the asymptotic normal interval is:
	\begin{align}
		E = Z_{\alpha/2} \sqrt{\frac{p(1-p)}{n}}.
	\end{align}
	\begin{enumerate}
		\item Using differential calculus or algebraic properties of quadratic functions, rigorously prove that the theoretical Bernoulli variance $g(p) = p(1-p)$ achieves an absolute global maximum on the domain $p \in [0,1]$ exactly at $p = 0.5$.
		\item Prove that if no prior information or reliable pilot studies on $p$ are available, the formula guaranteeing that the margin of error will never exceed the limit $E$ is:
		\begin{align}
			n = \left\lceil \frac{Z_{\alpha/2}^2}{4 E^2} \right\rceil.
		\end{align}
		\item A data science consulting firm needs to estimate voter intention regarding a public health initiative with a maximum margin of error of $3\%$ ($E = 0.03$) and a 95\% confidence level. Calculate the minimum number of valid surveys ($n$) that must be guaranteed.
	\end{enumerate}
\end{problema}

\begin{problema}
	\label{prob-ic-desaf-epistemologia}
	\textbf{Epistemological foundations and mathematical duality of the confidence interval:}
	Consider the formal construction of the empirical $(1-\alpha)$ confidence interval under normality:
	\begin{align}
		[L(X), \, U(X)] = \left[ \bar{X} - Z_{\alpha/2}\frac{\sigma}{\sqrt{n}}, \; \bar{X} + Z_{\alpha/2}\frac{\sigma}{\sqrt{n}} \right].
	\end{align}
	\begin{enumerate}
		\item Probabilistically prove prior to data collection (ex-ante) that the interval limits are random variables that exactly satisfy:
		\begin{align}
			P\left( L(X) \le \mu \le U(X) \right) = 1 - \alpha.
		\end{align}
		\item Suppose that once the particular real sample of 36 students from Problem \ref{prob-ic-1} is taken, the data are evaluated and the closed numerical interval $[74.08, \, 81.92]$ is published. Analytically explain why, from the strict axiomatic framework of Kolmogorov's frequentist statistics, it is theoretically incorrect to state that:
		\begin{quote}
			``The probability that the true mean score $\mu$ lies inside the interval $[74.08, \, 81.92]$ is $95\%$ (`$P(74.08 \le \mu \le 81.92) = 0.95$')''.
		\end{quote}
		\item Rigorously define what a ``95\% confidence'' means from the perspective of infinite repeated sampling over time, and contrast it with the epistemological concept of a \textbf{Bayesian credible interval}.
	\end{enumerate}
\end{problema}


\subsection{Detailed Solutions}

\subsubsection*{1. Fundamental Level (Simple / Direct)}

\begin{solucion}
	\begin{enumerate}
		\item \textbf{95\% confidence interval ($\alpha = 0.05$):}
		The critical quantile in the standardized normal distribution is $Z_{0.025} = 1.96$.
		The general normal interval formula is:
		\begin{align}
			IC_{95\%} = \bar{X} \pm Z_{\alpha/2} \frac{\sigma}{\sqrt{n}} = 78 \pm 1.96 \left( \frac{12}{\sqrt{36}} \right).
		\end{align}
		Simplifying the standard error of the mean ($\text{SE} = 12 / 6 = 2$):
		\begin{align}
			IC_{95\%} = 78 \pm 1.96(2) = 78 \pm 3.92 = [74.08, \; 81.92].
		\end{align}
		
		\item \textbf{99\% confidence interval ($\alpha = 0.01$):}
		For a 99\% confidence level, the critical value in each tail is $Z_{0.005} = 2.576$:
		\begin{align}
			IC_{99\%} = 78 \pm 2.576(2) = 78 \pm 5.152 = [72.85, \; 83.15].
		\end{align}
		
		\item \textbf{Comparison and theoretical justification of the width:}
		We calculate the width ($A = U - L$) for both intervals:
		\begin{align}
			A_{95\%} &= 81.92 - 74.08 = 7.84 \text{ points.} \\
			A_{99\%} &= 83.15 - 72.85 = 10.30 \text{ points.}
		\end{align}
		The 99\% interval has a substantially larger width ($10.30 > 7.84$). This occurs due to the uncertainty principle of frequentist inference: to increase our certainty (confidence) of not leaving the true parameter $\mu$ outside the interval when the amount of observed information ($n$) remains constant, we are forced to widen the margin of error, encompassing a larger fraction of the area under the tails of the normal bell curve.
	\end{enumerate}
\end{solucion}

\begin{solucion}
	\begin{enumerate}
		\item You must use \textbf{Student's $t$-distribution}. This decision is justified because two indisputable theoretical conditions hold:
		\begin{itemize}
			\item The sample size is small ($n = 16 < 30$), which prevents invoking the asymptotic normal approximation of the Central Limit Theorem.
			\item The population standard deviation $\sigma$ is unknown and must be replaced by the empirical sample standard deviation $s = 80$ hours, which introduces additional estimation error that requires the heavier tails of the $t$-distribution.
		\end{itemize}
		
		\item \textbf{Calculation of the 95\% interval with $\nu = 15$ degrees of freedom:}
		The exact degrees of freedom are $\nu = n - 1 = 16 - 1 = 15$. For $\alpha = 0.05$, the critical value from the Student's table is $t_{0.025, \, 15} = 2.131$.
		\begin{align}
			IC_{95\%} &= \bar{X} \pm t_{\alpha/2, \, \nu} \frac{s}{\sqrt{n}} = 1,200 \pm 2.131 \left( \frac{80}{\sqrt{16}} \right) \\
			&= 1,200 \pm 2.131 \left( \frac{80}{4} \right) = 1,200 \pm 2.131(20) = 1,200 \pm 42.62 \\
			&= [1,157.38, \; 1,242.62] \text{ hours.}
		\end{align}
		
		\item \textbf{Interpretation:} If we were to repeat this sampling process (drawing 16 lightbulbs at random and computing the interval) indefinitely many times, in $95\%$ of those samples the true average lifespan $\mu$ of the manufacturer's entire production would be captured between 1,157.38 and 1,242.62 operating hours.
	\end{enumerate}
\end{solucion}

\begin{solucion}
	\begin{enumerate}
		\item \textbf{Calculation of tail areas:}
		The $98\%$ confidence level corresponds to a central probability of $1 - \alpha = 0.98$.
		The total excluded complement or error risk is therefore:
		\begin{align}
			\alpha = 1 - 0.98 = 0.02 \text{ (2\%).}
		\end{align}
		By the even symmetry of the standardized normal function around zero, this outer mass splits equally between both tails of the distribution:
		\begin{align}
			\frac{\alpha}{2} = \frac{0.02}{2} = 0.01 \text{ in each extreme tail.}
		\end{align}
		
		\item \textbf{Exact critical value $Z_{\alpha/2}$:}
		The positive critical value $Z_{0.01}$ is the abscissa that accumulates to its left $(1 - 0.01) = 0.99$ of the total probability under the standard normal:
		\begin{align}
			\Phi(Z_{0.01}) = 0.9900.
		\end{align}
		Consulting inverse tables or using standard numerical probit interpolation, we obtain:
		\begin{align}
			Z_{0.01} = 2.326.
		\end{align}
	\end{enumerate}
\end{solucion}

\subsubsection*{2. Operational Level (Intermediate / Developmental)}

\begin{solucion}
	\begin{enumerate}
		\item \textbf{90\% confidence interval ($\alpha = 0.10$):}
		With $n = 25$, the degrees of freedom are $\nu = 24$. For 0.90 confidence, the table quantile is $t_{0.05, \, 24} = 1.711$:
		\begin{align}
			IC_{90\%} &= \bar{X} \pm t_{0.05, \, 24} \frac{s}{\sqrt{n}} = 15.0 \pm 1.711 \left( \frac{4.0}{\sqrt{25}} \right) \\
			&= 15.0 \pm 1.711(0.8) = 15.0 \pm 1.369 = [13.63, \; 16.37] \text{ hours.}
		\end{align}
		
		\item \textbf{95\% confidence interval ($\alpha = 0.05$):}
		For 0.95 confidence with $\nu = 24$, the critical value increases to $t_{0.025, \, 24} = 2.064$:
		\begin{align}
			IC_{95\%} &= 15.0 \pm 2.064(0.8) = 15.0 \pm 1.651 = [13.35, \; 16.65] \text{ hours.}
		\end{align}
		
		\item \textbf{Methodologies to improve precision and narrow the 95\% interval:}
		To reduce the width of the interval while maintaining the 95\% confidence level ($Z$ or $t$ fixed), the researcher has two mathematical avenues:
		\begin{itemize}
			\item \textbf{Increase the sample size ($n$):} Because the standard error of the estimate is inversely proportional to the square root of $n$ ($\text{SE} = s / \sqrt{n}$), quadrupling the sample size (e.g., increasing from $n=25$ to $n=100$) halves the margin of error uncertainty, producing a much tighter interval.
			\item \textbf{Improve the sampling design or reduce empirical variability ($s$):} Stratifying the sample by academic faculty or field of study, or employing standardized automated time-tracking apps to avoid memory recall bias, can reduce internal dispersion in the observed data ($s$), making the interval more compact.
		\end{itemize}
	\end{enumerate}
\end{solucion}

\begin{solucion}
	\begin{enumerate}
		\item \textbf{Construction of the 95\% interval:}
		Since this is a large sample ($n = 100 \ge 30$), the Central Limit Theorem guarantees that the sample mean is asymptotically normally distributed. With $Z_{0.025} = 1.96$:
		\begin{align}
			IC_{95\%} &= \bar{X} \pm Z_{0.025} \frac{s}{\sqrt{n}} = 4.5 \pm 1.96 \left( \frac{2.0}{\sqrt{100}} \right) \\
			&= 4.5 \pm 1.96 \left( \frac{2.0}{10} \right) = 4.5 \pm 1.96(0.2) = 4.5 \pm 0.392 \\
			&= [4.108, \; 4.892] \text{ minutes.}
		\end{align}
		
		\item \textbf{Estimation margin of error ($E$):}
		The half-width of uncertainty of the current interval is $E = 0.392$ minutes (approximately 23.5 seconds).
		
		\item \textbf{Derivation of the new sample size for a preset bound $E = 0.2$:}
		We start from the algebraic definition of the margin of error:
		\begin{align}
			E = Z_{\alpha/2} \frac{s}{\sqrt{n}}.
		\end{align}
		Solving for the square root variable $\sqrt{n}$:
		\begin{align}
			\sqrt{n} = \frac{Z_{\alpha/2} \cdot s}{E}.
		\end{align}
		Squaring both sides yields the required sample size formula:
		\begin{align}
			n = \left( \frac{Z_{\alpha/2} \cdot s}{E} \right)^2.
		\end{align}
		Substituting the study values ($Z_{0.025} = 1.96$, $s = 2.0$, $E = 0.2$):
		\begin{align}
			n = \left( \frac{1.96 \times 2.0}{0.2} \right)^2 = \left( \frac{3.92}{0.2} \right)^2 = (19.6)^2 = 384.16.
		\end{align}
		Since we cannot sample a fraction of a call, we apply the ceiling function ($\lceil x \rceil$), always rounding up to the next integer:
		\begin{align}
			n_{\text{optimal}} = \lceil 384.16 \rceil = 385 \text{ calls.}
		\end{align}
		The company must collect a sample of at least 385 total calls to guarantee a margin of error equal to or less than 12 seconds ($0.2$ minutes).
	\end{enumerate}
\end{solucion}

\begin{solucion}
	\begin{enumerate}
		\item \textbf{Construction of the 95\% interval:}
		For $n = 25$ batteries, we apply Student's $t$-distribution with $\nu = 25 - 1 = 24$ degrees of freedom. The associated critical quantile is $t_{0.025, \, 24} = 2.064$:
		\begin{align}
			IC_{95\%} &= \bar{X} \pm t_{0.025, \, 24} \frac{s}{\sqrt{n}} = 490 \pm 2.064 \left( \frac{30}{\sqrt{25}} \right) \\
			&= 490 \pm 2.064 \left( \frac{30}{5} \right) = 490 \pm 2.064(6) = 490 \pm 12.384 \\
			&= [477.62, \; 502.38] \text{ recharge cycles.}
		\end{align}
		
		\item \textbf{Evaluation of the advertising threshold:}
		The value claimed by the advertising company is $\mu_0 = 500$ cycles. We numerically observe that $477.62 < 500 < 502.38$, meaning that the advertised threshold \textbf{is indeed contained} within the confidence interval.
		
		\item \textbf{Technical conclusion and equivalence with hypothesis testing:}
		Although the experimental sample yielded a mean below the specification ($490 < 500$), the standard error of the experiment ($6$ cycles) multiplied by Student's quantile proves that the observed discrepancy of $-10$ cycles is completely admissible as a typical random sampling fluctuation.
		
		By the fundamental theoretical duality between interval estimation and hypothesis testing, since the hypothesized value $\mu_0 = 500$ falls inside the two-tailed $(1-\alpha) = 0.95$ confidence interval, if we performed a formal hypothesis test $H_0: \mu = 500$ vs. $H_a: \mu \neq 500$ at significance level $\alpha = 0.05$, \textbf{there would not be sufficient statistical evidence to reject the null hypothesis}. Consequently, the laboratory concludes that the manufacturer's advertising claim cannot be conclusively refuted.
	\end{enumerate}
\end{solucion}

\subsubsection*{3. Analytical Level (Advanced / Connection)}

\begin{solucion}
	\begin{enumerate}
		\item \textbf{Calculation of the standard error of the difference between independent means without assuming homogeneous variances ($\sigma_1^2 \neq \sigma_2^2$):}
		By independence of the two study groups, the variance of the difference between sample means is the sum of the individual variances ($\Var(\bar{X}_1 - \bar{X}_2) = \Var(\bar{X}_1) + \Var(\bar{X}_2)$). Its square root is the estimated standard error:
		\begin{align}
			\text{SE}(\bar{X}_1 - \bar{X}_2) = \sqrt{\frac{s_1^2}{n_1} + \frac{s_2^2}{n_2}}.
		\end{align}
		Substituting the observed sample variances for men ($s_1 = 300, n_1 = 40$) and women ($s_2 = 250, n_2 = 45$):
		\begin{align}
			\text{SE}(\bar{X}_1 - \bar{X}_2) &= \sqrt{\frac{300^2}{40} + \frac{250^2}{45}} = \sqrt{\frac{90,000}{40} + \frac{62,500}{45}} \\
			&= \sqrt{2,250 + 1,388.889} = \sqrt{3,638.889} \approx 60.323 \text{ kcal/day.}
		\end{align}
		
		\item \textbf{Construction of the 95\% interval for $\mu_1 - \mu_2$:}
		The basic difference between sample averages is:
		\begin{align}
			\hat{D} = \bar{X}_1 - \bar{X}_2 = 2,500 - 2,100 = 400 \text{ kcal/day.}
		\end{align}
		Because both samples comfortably exceed 30 observations ($n_1 = 40, n_2 = 45$), we invoke the asymptotic normal approximation of the Central Limit Theorem with $Z_{0.025} = 1.96$:
		\begin{align}
			IC_{95\%}(\mu_1 - \mu_2) &= (\bar{X}_1 - \bar{X}_2) \pm Z_{0.025} \cdot \text{SE}(\bar{X}_1 - \bar{X}_2) \\
			&= 400 \pm 1.96(60.323) = 400 \pm 118.233 \\
			&= [281.77, \; 518.23] \text{ kcal/day.}
		\end{align}
		
		\item \textbf{Biological statistical significance conclusion:}
		The confidence interval obtained for the caloric gap is bounded between $281.77$ kcal and $518.23$ kcal. Since the lower bound is strictly positive ($281.77 > 0$), \textbf{the value zero is not included inside the confidence interval}.
		
		This provides compelling statistical evidence to affirm, with 95\% confidence, that adult men in this target population consume a significantly higher amount of daily calories than women, exhibiting a biological difference that ranges between approximately 282 and 518 extra kilocalories per day on average.
	\end{enumerate}
\end{solucion}

\begin{solucion}
	\begin{enumerate}
		\item \textbf{Derivation and calculation of the pooled sample variance ($S_p^2$):}
		When assuming that the two normal populations share an identical common variance ($\sigma_1^2 = \sigma_2^2 = \sigma^2$), the individual estimates $s_1^2$ and $s_2^2$ are two independent observations estimating the exact same parameter $\sigma^2$. To maximize degrees of freedom without introducing bias, we construct the \textbf{pooled variance ($S_p^2$)}, which is the weighted average of the sample variances weighted by their respective internal degrees of freedom ($n_i - 1$):
		\begin{align}
			S_p^2 = \frac{(n_1 - 1)s_1^2 + (n_2 - 1)s_2^2}{n_1 + n_2 - 2}.
		\end{align}
		Substituting the data for Catalysts A ($n_1=10, s_1=4.0$) and B ($n_2=12, s_2=5.0$):
		\begin{align}
			S_p^2 &= \frac{(10 - 1)(4.0)^2 + (12 - 1)(5.0)^2}{10 + 12 - 2} = \frac{9(16.0) + 11(25.0)}{20} \\
			&= \frac{144.0 + 275.0}{20} = \frac{419.0}{20} = 20.95 \implies S_p = \sqrt{20.95} \approx 4.5771\%.
		\end{align}
		
		\item \textbf{Construction of the difference interval with $\nu = 20$:}
		The standard error of the difference using pooled variance is:
		\begin{align}
			\text{SE}_p(\bar{X}_1 - \bar{X}_2) = \sqrt{S_p^2 \left( \frac{1}{n_1} + \frac{1}{n_2} \right)} = \sqrt{20.95 \left( \frac{1}{10} + \frac{1}{12} \right)} = \sqrt{20.95(0.18333)} \approx 1.9598\%.
		\end{align}
		The observed difference between yields is $\hat{D} = 85.0\% - 81.0\% = 4.0\%$.
		For exact degrees of freedom $\nu = 10 + 12 - 2 = 20$ and a 95\% confidence level, the value in the Student tables is $t_{0.025, \, 20} = 2.086$. Therefore:
		\begin{align}
			IC_{95\%}(\mu_A - \mu_B) &= (\bar{X}_1 - \bar{X}_2) \pm t_{0.025, \, 20} \cdot \text{SE}_p \\
			&= 4.0 \pm 2.086(1.9598) = 4.0 \pm 4.088 \\
			&= [-0.088\%, \; +8.088\%].
		\end{align}
		\textbf{Interpretation of experimental result:} Although the nominal yield of Catalyst A exceeds that of Catalyst B by $4\%$, the confidence interval spans the value zero (ranging from $-0.09\%$ up to $+8.09\%$). Consequently, at 95\% confidence in these small samples, Catalyst A cannot formally and statistically be declared superior to Catalyst B.
	\end{enumerate}
\end{solucion}

\subsubsection*{4. Challenging Level (Experts / Deep Deduction)}

\begin{solucion}
	\begin{enumerate}
		\item \textbf{Analytical maximization of the elementary Bernoulli variance:}
		Let us define the parametric variance function on the unit interval:
		\begin{align}
			g(p) = p(1 - p) = p - p^2, \quad p \in [0, 1].
		\end{align}
		We differentiate the function with respect to the proportion $p$ and set it equal to zero to find critical points:
		\begin{align}
			g'(p) = \frac{d}{dp}(p - p^2) = 1 - 2p = 0 \implies 2p = 1 \implies p^* = 0.5.
		\end{align}
		We check the sufficient condition for an absolute maximum via the second derivative:
		\begin{align}
			g''(p) = \frac{d^2}{dp^2}(p - p^2) = -2 < 0.
		\end{align}
		Because the concavity is uniformly negative throughout its entire domain, it is irrevocably proven that $g(p)$ reaches an absolute global supremum right when the event exhibits maximum uncertainty ($p = 0.5$). The numerical value of this upper bound is:
		\begin{align}
			g(0.5) = 0.5(1 - 0.5) = 0.5 \times 0.5 = 0.25 = \frac{1}{4}.
		\end{align}
		
		\item \textbf{Rigorously deriving the sample size bound:}
		To guarantee that a $(1-\alpha)$ confidence interval has a half-width less than or equal to a preset margin $E$, we impose the general algebraic condition:
		\begin{align}
			Z_{\alpha/2} \sqrt{\frac{p(1-p)}{n}} \le E.
		\end{align}
		Squaring both sides of the inequality:
		\begin{align}
			Z_{\alpha/2}^2 \frac{p(1-p)}{n} \le E^2 \implies n \ge \frac{Z_{\alpha/2}^2 \cdot p(1-p)}{E^2}.
		\end{align}
		Since the true population proportion $p$ is unknown a priori by the surveyor, for the inequality $n \ge \dots$ to hold true \textbf{regardless of what true value $p$ takes in reality}, we must replace the variable numerator with its extreme uniform upper bound proven in Part 1 ($p(1-p) \le 0.25$):
		\begin{align}
			n \ge \frac{Z_{\alpha/2}^2 (0.25)}{E^2} = \frac{Z_{\alpha/2}^2}{4 E^2}.
		\end{align}
		This yields the famous conservative sample size formula for proportions:
		\begin{align}
			n = \left\lceil \frac{Z_{\alpha/2}^2}{4 E^2} \right\rceil.
		\end{align}
		
		\item \textbf{Applied calculation for the consulting firm ($E = 0.03$, 95\% confidence):}
		We substitute $Z_{0.025} = 1.96$ and $E = 0.03$ into the derived equation:
		\begin{align}
			n &= \frac{(1.96)^2}{4 (0.03)^2} = \frac{3.8416}{4 (0.0009)} = \frac{3.8416}{0.0036} = 1,067.111.
		\end{align}
		Applying the ceiling function to round up to exact integers so as not to sacrifice even a thousandth of precision in the final margin:
		\begin{align}
			n_{\text{optimal}} = \lceil 1,067.111 \rceil = 1,068 \text{ valid surveys.}
		\end{align}
	\end{enumerate}
\end{solucion}

\begin{solucion}
	\begin{enumerate}
		\item \textbf{Ex-ante probabilistic proof (prior to sampling):}
		Let $\bar{X} \sim N\left(\mu, \dfrac{\sigma^2}{n}\right)$ be the random variable describing the sample average prior to collecting data. By Central Limit theory, the standardized variable satisfies:
		\begin{align}
			Z = \frac{\bar{X} - \mu}{\sigma/\sqrt{n}} \sim N(0,1).
		\end{align}
		By definition of symmetric quantiles at level $(1-\alpha)$, the random event $Z$ is contained in the interval $[-Z_{\alpha/2}, \, +Z_{\alpha/2}]$ with exact probability:
		\begin{align}
			P\left( -Z_{\alpha/2} \le \frac{\bar{X} - \mu}{\sigma/\sqrt{n}} \le +Z_{\alpha/2} \right) = 1 - \alpha.
		\end{align}
		We multiply the entire double inequality inside the probability operator by the strictly positive standard error scale constant $\dfrac{\sigma}{\sqrt{n}}$:
		\begin{align}
			P\left( -Z_{\alpha/2}\frac{\sigma}{\sqrt{n}} \le \bar{X} - \mu \le +Z_{\alpha/2}\frac{\sigma}{\sqrt{n}} \right) = 1 - \alpha.
		\end{align}
		We subtract the random sample mean $\bar{X}$ across the three interior terms:
		\begin{align}
			P\left( -\bar{X} - Z_{\alpha/2}\frac{\sigma}{\sqrt{n}} \le -\mu \le -\bar{X} + Z_{\alpha/2}\frac{\sigma}{\sqrt{n}} \right) = 1 - \alpha.
		\end{align}
		We multiply by the negative constant $(-1)$ across the entire inequality, recalling that in algebra this reverses the direction of the inequalities:
		\begin{align}
			P\left( \bar{X} + Z_{\alpha/2}\frac{\sigma}{\sqrt{n}} \ge \mu \ge \bar{X} - Z_{\alpha/2}\frac{\sigma}{\sqrt{n}} \right) = 1 - \alpha.
		\end{align}
		Rearranging the bounds from lowest to highest, we define the lower bound function $L(X) = \bar{X} - Z_{\alpha/2}\frac{\sigma}{\sqrt{n}}$ and upper bound function $U(X) = \bar{X} + Z_{\alpha/2}\frac{\sigma}{\sqrt{n}}$, fully proving the theorem:
		\begin{align}
			P\left( L(X) \le \mu \le U(X) \right) = 1 - \alpha.
		\end{align}
		
		\item \textbf{Epistemological failure of the ex-post probabilistic statement:}
		In the Classical Frequentist statistical paradigm formulated by Neyman and Pearson, the true population parameter $\mu$ is a \textbf{fixed, immutable, and unique constant of nature}, completely devoid of randomness (it is not a random variable).
		
		Once a sample is drawn from the physical world and the numbers are plugged into the formula, the random variable $\bar{X}$ collapses into the constant observed realization $\bar{x} = 78$, turning the interval into the rigidly fixed numerical set $[74.08, \, 81.92]$.
		
		At this ex-post stage, \textbf{all elements inside the expression are deterministic constants}. Therefore, the event $\mu \in [74.08, \, 81.92]$ has no associated randomness whatsoever. To an omniscient observer, only two mathematically valid states exist for that particular probability:
		\begin{align}
			P(74.08 \le \mu \le 81.92) = 
			\begin{cases}
				1, & \text{if the true } \mu \text{ actually lies inside this range.} \\
				0, & \text{if the true } \mu \text{ fell outside this realization.}
			\end{cases}
		\end{align}
		Assigning a decimal or percentage number like ``$0.95$'' to that specific empirical interval is a formal conceptual fallacy, mistaking the probability of the sampling procedure for the probability of a fixed parameter.
		
		\item \textbf{Correct axiomatic definition and contrast with the Bayesian paradigm:}
		\begin{itemize}
			\item \textbf{Frequentist Interpretation (95\% Confidence):} Confidence is an exclusive property of the \textbf{interval-generating procedure or process}, not of any single realized numerical interval. It means that if we repeated the random sampling experiment an infinite number of times, generating thousands of intervals $[L(X_k), U(X_k)]$ with the different averages arising by chance, exactly \textbf{$95\%$ of all random intervals constructed by the formula will successfully enclose the true parameter $\mu$}, while the remaining $5\%$ will miss due to marginal sampling fluctuation.
			\item \textbf{Contrast with the Bayesian Credible Interval:} If a scientist or data scientist truly needs to make direct, intuitive probabilistic statements about where the unknown parameter lies once data are observed, they must depart from the frequentist paradigm and adopt \textbf{Bayesian analysis}. In Bayesian inference, the parameter $\theta$ is explicitly treated as a random variable endowed with a prior distribution $\pi(\theta)$. Conditioning on the observed sample data $X=x$, Bayes' Theorem computes the \textbf{posterior distribution} $\pi(\theta \mid X=x)$. From this distribution, one constructs the \emph{Bayesian Credible Interval} $[A, B]$ such that:
			\begin{align}
				P(A \le \theta \le B \mid X = x) = 0.95.
			\end{align}
			Here, it is rigorous, valid, and epistemologically sound to state that, given the data and prior information, there is a $95\%$ probability that the parameter $\theta$ currently resides within the calculated boundaries.
		\end{itemize}
	\end{enumerate}
\end{solucion}
